\section{Replacing branch disks by trivial covers}
\label{sec:patching}

We first isolate the analytic operation used in both stages of the proof.
It replaces a branched cover by an étale cover after another cover supplies
roots of the boundary coordinates.  The roots must exist over the entire
inverse image of each boundary annulus.  The resulting analytic map is
needed only outside the branch disks.
This is an application of the patching method for covers developed in
formal geometry by Harbater, in rigid geometry by Liu, and in the
Berkovich setting by Poineau
\cite{HarbaterPatching,LiuPatching,GAGA}.  Here the root on every inverse
collar supplies the trivialization needed to replace the whole branched
cover over a disk, including all of its sheets.

\begin{lemma}\label{lem:local-normal-form}
Let $f:C\to X_{\bar k}$ be a finite surjective morphism of smooth proper
connected curves, and let $b\in X(\bar k)$.  Fix a rational function $t_b$
which is a local parameter at $b$.  There is an open disk $B_b$ about $b$
in $X_\Omega^{\mathrm{an}}$, with coordinate $t_b$, such that its full
inverse image in $C_\Omega^{\mathrm{an}}$ is a disjoint union of open
disks.  On the disk about $c\in f^{-1}(b)$ there is a coordinate $z_c$
for which
\[
                         t_b=z_c^{e_c},
\]
where $e_c$ is the ramification index at $c$.  No condition that $e_c$ be
prime to $p$ is required.  For finitely many distinct points $b$, the disks
can be chosen disjoint and disjoint from the fixed point $\xi$.
\end{lemma}

\begin{proof}
Choose a local analytic parameter $u$ at $c$.  The pullback of $t_b$ is
$u^{e_c}a(u)$ with $a(0)\ne0$.  After choosing an $e_c$-th root of
$a(0)$ in $\Omega$ and shrinking the disk, $a(u)/a(0)$ is sufficiently
close to $1$ for its $e_c$-th root to be given by
\[
 \exp\!\left(e_c^{-1}\log(a(u)/a(0))\right).
\]
Here it is enough to make the valuation of $a(u)/a(0)-1$ greater than
$v_p(e_c)+1/(p-1)$.  Thus $z_c=u\,a(u)^{1/e_c}$ is analytic with a
nonzero derivative at the center.  Analytic inversion gives a smaller
disk on which it is a coordinate and $t_b=z_c^{e_c}$.

Do this in disjoint neighborhoods of all the finitely many points above
$b$.  The complement of these neighborhoods in the proper analytic curve
$C_\Omega^{\mathrm{an}}$ is compact.  Its image under $f$ is compact,
closed, and does not contain $b$.  Consequently a sufficiently small disk
about $b$ has its entire inverse image in the chosen neighborhoods.
Shrink it once more so that, in every $z_c$-chart, that inverse image is
the disk defined by the corresponding bound on $|z_c|^{e_c}$.

The point $\xi$ does not equal any point of $X(\bar k)$, since its
projection to $X_{\bar k}$ is generic.  The shrinking disks at $b$ form
a neighborhood basis there.  Hausdorffness therefore permits all disks
to avoid $\xi$ and one another.  We may choose their radii in
$|\Omega^\times|$.  The argument also applies when $\xi$ is an
$\Omega$-rational point whose projection to $X_{\bar k}$ is generic.
\end{proof}

\begin{proposition}\label{prop:patching}
Let $f:C\to X_{\bar k}$ be as in Lemma~\ref{lem:local-normal-form},
of degree $d$, with branch set $B$.  Choose pairwise disjoint disks $B_b$
as in that lemma.
For each $b\in B$, let $m_b$ be a positive multiple of all ramification
indices above $b$.  Suppose that $g:Y\to X_\Omega$ is a finite
surjective morphism from a smooth proper connected curve and that there
are radii $0<r_b<R_b$ in $|\Omega^\times|$, strictly inside $B_b$, such
that $g^*t_b$ has an analytic $m_b$-th root on the whole inverse image of
\[
                       A_b=\{r_b<|t_b|<R_b\}.
\]
Set
\[
 U=X_\Omega^{\mathrm{an}}\setminus
       \bigcup_{b\in B}\{|t_b|\le r_b\},
\]
where each indicated closed disk is taken inside $B_b$.
Then there is a smooth proper connected curve $V_\Omega$ and a finite
étale surjective map $V_\Omega\to Y$, together with an analytic map
\[
 (V_\Omega^{\mathrm{an}})|_U\longrightarrow C_\Omega^{\mathrm{an}}
 \quad\text{over }X_\Omega^{\mathrm{an}}.
\]
The notation on the left denotes the inverse image of $U$ under the
composite map to $X_\Omega^{\mathrm{an}}$.
\end{proposition}

\begin{proof}
The inverse image $Y_U$ of $U$ and the inverse images $Y_b$ of the disks
$\{|t_b|<R_b\}$ form an open cover of $Y^{\mathrm{an}}$.  Distinct
$Y_b$ are disjoint, and $Y_U\cap Y_b=g^{-1}(A_b)$.

On $Y_U$, take the pullback of $C_\Omega^{\mathrm{an}}\to
X_\Omega^{\mathrm{an}}$.  It is finite étale of degree $d$, since $U$
contains no branch point.  On each $Y_b$, take the trivial degree-$d$
cover.  These two covers are isomorphic over their overlap.  Indeed,
write $q_b^{m_b}=g^*t_b$ for the assumed analytic root.  In the local
factor corresponding to $c\in f^{-1}(b)$, put
$a_c=q_b^{m_b/e_c}$.  The polynomial defining that factor splits as
\[
 Z^{e_c}-g^*t_b=\prod_{\zeta\in\mu_{e_c}}(Z-\zeta a_c).
\]
Each $a_c$ is an analytic unit on the overlap.  Distinct linear factors
are comaximal, because $(\zeta-\zeta')a_c$ is a unit.  This is true also
when $p\mid e_c$: a nonzero constant in $\Omega$ is invertible, regardless
of its absolute value.  Taking all factors in
Lemma~\ref{lem:local-normal-form} gives a trivialization of the full
degree-$d$ cover, not just of one sheet.

Choose these trivializations and glue the finite locally free analytic
algebras.  There is no further cocycle condition between different branch
disks, as those disks are disjoint.  The glued algebra is coherent and
finite étale over $Y^{\mathrm{an}}$, these being local properties.
Proper analytic GAGA \cite[Appendix~A, Theorem~A.1(ii)]{GAGA} algebraizes the coherent algebra and
its unit and multiplication; full faithfulness preserves the algebra
identities.  Equivalently, one uses the corresponding equivalence for
finite covers of a proper curve.  \'Etaleness is detected after
analytification \cite[Fact~5.1.5]{TemkinAnalytic}, so its relative spectrum is a finite étale cover of $Y$.

Choose any nonempty connected component and call it $V_\Omega$.
A connected component of this finite étale cover has positive locally
constant degree over the connected curve $Y$, hence is still surjective.
It is smooth and proper.  Restricting the original projection from the
pullback over $Y_U$ gives the asserted analytic map.  If $B$ is empty,
the whole pullback is already finite étale and the same component
argument applies without gluing.
\end{proof}

We record explicitly how this analytic map produces an embedding into
the fixed completion $H$.  This step concerns the entire $H$-valued
point, not just its underlying analytic point.

\begin{lemma}\label{lem:point-transfer}
In Proposition~\ref{prop:patching}, suppose that the composite
$V_\Omega\to X_\Omega$ is isomorphic, over $X_\Omega$, to the base
change of a finite morphism $V\to X_{\bar k}$ from a smooth proper
connected curve, and assume that $\xi\in U$.
If $\bar k(V)$ admits a $K_0$-embedding into $H$, then
$\bar k(C)$ admits a $K_0$-embedding into $H$.
\end{lemma}

\begin{proof}
The given embedding determines an $H$-valued point of $V$.  Together
with the fixed inclusion $\Omega\subset H$, it determines an
$H$-valued point of $V\times_{\bar k}\Omega$, and the chosen
isomorphism transports this point to $V_\Omega$.  Its image in
$X_\Omega(H)$ is exactly the fixed point induced by $K_0\hookrightarrow
H$ and $\Omega\hookrightarrow H$.  In particular its underlying
analytic image is $\xi\in U$.

An $H$-valued point of an algebraic variety over $\Omega$ is equivalently
an $H$-valued point of its analytification \cite[Fact~5.1.3]{TemkinAnalytic}.
More explicitly, it gives an
analytic point $v$ together with an isometric embedding of the completed
residue field $\mathcal H(v)$ into $H$.  The analytic map over $U$ in
Proposition~\ref{prop:patching} induces a compatible map on completed
residue fields, and hence an $H$-valued point of $C_\Omega$ over the
same point of $X_\Omega(H)$.  Project this point to $C$.  Its image on
$X_{\bar k}$ is the generic point with the specified map $K_0\to H$.
Since $C\to X_{\bar k}$ is finite, its support on $C$ is also generic.
The point therefore induces the required field embedding
$\bar k(C)\hookrightarrow H$ over $K_0$.

Neither the analytic map over $U$ nor a chosen local trivialization has
been asserted to descend to $\bar k$.
\end{proof}

We shall form auxiliary covers by pullback and compositum, then pass to
a finite étale upper cover.  The following lemma shows that these
operations preserve the bound on horizontal ramification needed for
membership in $L_\ell$.

\begin{lemma}\label{lem:horizontal-closure}
Fix a positive integer $n$.  Among finite extensions of $K_0$, the
condition that every ramification index at a closed point of
$X_{\bar k}$ divide $n$ is preserved by taking subextensions, finite
composita, and Galois closures.  It is also preserved by normalized
pullback along a finite morphism of smooth proper curves and by
composition on top with a finite étale cover.
Consequently every finite cover of $X_{\bar k}$ with indices dividing
$\ell$ has a $K_0$-embedding of its function field into $L_\ell$.
\end{lemma}

\begin{proof}
All valuations in this statement are the horizontal, discrete
closed-point valuations, which are trivial on $\bar k$.  Their residue
characteristic is zero.  The completed base field at such a point is
$A=\bar k((t))$.  Every finite extension of $A$ has residue field
$\bar k$ and is of the form $\bar k((u))$ for a uniformizer $u$:
successive subtraction of residue representatives gives the power-series
expansion in its complete valuation ring.  If its ramification index is
$e$, then $t=u^e a(u)$ for a unit $a(u)\in\bar k[[u]]^\times$.
The unit has an $e$-th root in $\bar k[[u]]$, by Hensel's lemma in
residue characteristic zero.  Changing the parameter gives
$t=z^e$.  Thus, inside a fixed algebraic closure, the local field is
$A(t^{1/e})$, which is cyclic over $A$ since $\bar k$ contains all roots
of unity.

For $e\mid n$, this is a subfield of $A(t^{1/n})$.  All local factors
of a finite compositum, or of a normal closure formed from conjugates,
therefore remain subfields of this same cyclic extension.  This proves
the assertions about subextensions, composita, and Galois closures.

For pullback, write the local base map as $t=s^r a(s)$, where $r\ge1$
and $a(s)$ is a unit.  Extracting its $e$-th root in $\bar k[[s]]$
shows that every normalized local factor after pullback has ramification
index $e/\gcd(e,r)$, which divides $e$.  Finally, ramification indices
multiply in a tower, and an étale upper cover contributes index $1$.
Apply the Galois-closure assertion to a cover with indices dividing
$\ell$ and embed that closure into the fixed $K^{\mathrm{sep}}$.
Its field lies in the defining compositum $L_\ell$.
\end{proof}

\begin{remark}
Lemma~\ref{lem:horizontal-closure} makes no assertion about tame
reduction in residue characteristic $p$, nor about composing two
arbitrarily ramified covers.  Its last assertion about composition
requires the upper cover to be étale.
\end{remark}
